\documentclass[a4paper,11pt]{article}
\usepackage[justification=centering]{caption}
\usepackage[justification=centering]{subcaption}
\usepackage[inkscapelatex=true,inkscapearea=page]{svg}
\usepackage{enumitem}
\usepackage{bm,amsmath,amssymb}

\usepackage{fontspec}
\defaultfontfeatures{Mapping=tex-text}
\setromanfont{Linux Libertine O}
\usepackage{libertine}

\usepackage[dvipsnames]{xcolor}
\usepackage{natbib}\bibliographystyle{authordate1}
\usepackage[a4paper,hmargin=2cm,vmargin=3cm]{geometry}
\usepackage[colorlinks=true,allcolors=MidnightBlue]{hyperref}
\usepackage{float}

\providecommand{\tightlist}{\setlength{\itemsep}{0pt}\setlength{\parskip}{0pt}}
\setlength{\parindent}{0pt}
\svgsetup{width=0.9\linewidth}

\title{Derivation of the double-averaged Navier-Stokes equations for a fixed porous medium}
\author{Axel Giboulot}
\date{\today}

\renewcommand*{\figureautorefname}{figure}
\renewcommand*{\tableautorefname}{table}
\renewcommand*{\equationautorefname}{equation}

\begin{document}

% \renameautorefname{Figure}{figure}
% \renameautorefname{Table}{table}
% \renameautorefname{Equation}{equation}

  \maketitle

{\hypersetup{hidelinks}\tableofcontents}

\clearpage

\section{Introduction}

The Navier-Stokes equations describe the conservation of mass and momentum of fluids in motion.
They can be averaged in both space and time to explicitly state the contributions of turbulence and the resistance to flow offered by a porous medium.
The resulting set of equations is called the \emph{double-averaged} Navier-Stokes equations.
It is available in literature \citep{Nikora2007}, but its derivation can be tedious.
While derivations have already been published \citep{Rousseau2022}, a few steps could benefit from further clarification.
Consequently, I hereby derive this set of equations starting from the well-known time-(or Reynolds-)averaged Navier-Stokes equations in an editable format. Contributions are welcome.

\section{The time-averaged Navier-Stokes equations}

Newton's second law can be applied to a Newtonian fluid to find the local formulation of momentum balance. Combined with the local mass conservation equation, they make the Navier-Stokes equations.
Then decomposing the velocity vector $\bm u$ into a time-averaged component $\bm{\bar u}$ and a temporal fluctuation $\bm u'$ yields the following time-(or Reynolds-)averaged Navier-Stokes equations. In this derivation, the fluid is considered incompressible.

$$\bm\nabla\cdot\bm{\bar u} = 0,$$

$$\frac{\partial\bm{\bar u}}{\partial t} + \bm\nabla\cdot\left(\bm{\bar u}\otimes \bm{\bar u}\right) + \bm\nabla\cdot\overline{\bm u'\otimes\bm u'} = \bm g -\frac{1}{\varrho}\nabla p + \nu\nabla^2\bm{\bar u}.$$

Where

\begin{itemize}
  \item $\varrho$ is the unit mass of the fluid,
  \item $\bm g$ stands for gravity,
  \item $p$ is the pressure,
  \item $\nu$ is the kinematic viscosity and
  \item $t$ is time.
\end{itemize}

The time-averaging procedure over some period $T$ is defined for any variable $\psi$ as follows with $t$ representing time and $\bm x$ the vector of spatial coordinates.

$$\bar \psi(\bm x, t)\doteq\frac{1}{T}\int_{t-T/2}^{t+T/2} \psi(\bm x, t)~\mathrm dt,$$
$$\psi = \bar \psi + \psi'.$$

It is assumed that the following time-averaging rule holds for all well-behaved functions $\omega$ and constants $a$ and $b$.

$$\overline{a\bar\psi\omega+b} = a\bar\psi\bar\omega+b.$$

It then comes naturally that the time-average of the temporal fluctuations $\overline{\psi'}$ is null.

$$\bar\psi=\overline{\bar\psi}+\overline{\psi'},$$
$$\Rightarrow \bar\psi=\bar\psi+\overline{\psi'},$$
$$\Rightarrow \overline{\psi'}=0.$$

\section{Space averaging}

We wish to spatially average the previous Reynolds-(or time-)averaged equations over a porous medium. This can be done with a decomposition like Reynolds' by defining a spatial average $\langle\bm{\bar u}\rangle_f$ over some control volume $\mathcal V$ and its underlying spatial fluctuations $\tilde{\bm u}$.

\begin{equation}\label{eq:decomposition}
  {\bm{\bar u} = \langle\bm{\bar u}\rangle_f + \tilde{\bm u}}
\end{equation}

This procedure allows to smoothen-out the micro-scale features of the system under investigation while retaining its macro-scale features. This scale of observation is known as the meso-scale.
The control volume should be smaller than the size of the system but much larger than the small features of the porous media.
This control volume is sometimes refered to as the representative elementary volume (REV) and is commonly used in geomechanics where the flow is usually laminar. The double-averaged Navier-Stokes equations extend this concept to turbulent flows.

\section{Definition of the averages}

To formally delimit the porous medium, the indicating function ${\delta(\bm x, t)}$ is defined : its value is 1 in the fluid phase $\mathcal V_f$ and 0 in the solid phase $\mathcal V_s$.
A drawing of the situation is provided in \autoref{fig:def}.

$$\delta(\bm x, t) \doteq
\begin{cases}
    1 \text{ if } \bm x \in \mathcal V_f,\\
    0 \text{ otherwise.}
\end{cases}$$

\begin{figure}
  \centering
  \includesvg{gfx/Whitaker_definition.svg}
  \caption{Drawing of an control volume \citep{Whitaker1999}.\label{fig:def}}
\end{figure}

Two kinds of averages are needed :

\begin{itemize}
\item
  a first average over the control volume with a fixed averaging window and
\item
  a second average over the fluid phase where all variables are clearly defined.
\end{itemize}

\subsection{The fluid phase average}

This average is defined as the integral of some quantity $\psi$ in the fluid phase, normalized with respect to the entire control volume $\mathcal V$.

$$\langle\psi\rangle_v(\bm x, t) \doteq \frac{1}{V}\int_{\mathcal V} \psi(\bm x+\bm\xi, t) \delta(\bm x + \bm\xi, t) ~\left|\mathrm d\bm \xi\right|.$$

Where the position $\bm x$ is fixed and the variable ${\bm x + \bm\xi}$ is the vector exploring the domain $\mathcal V$.
The quantity ${\left|\mathrm d\bm\xi\right|}$ is defined as the associated infinitesimal volume.
This formulation can be simplified by noting that only the subdomain $\mathcal V_f\subseteq\mathcal V$ contributes to the integral, as the indicating function $\delta$ nullifies the product elsewhere (in ${\mathcal V_s=\mathcal V\backslash\mathcal V_f}$).

$$\langle\psi\rangle_v(\bm x, t) = \frac{1}{V}\int_{\mathcal V_f} \psi(\bm x+\bm\xi, t)~\left|\mathrm d\bm \xi\right|.$$

For the sake of readability, the notation is further simplified by varying the position $\bm x$ in the right-hand side over the subdomain $\mathcal V_f$ with the control volume $\mathcal V$ being centered on the position $\bm x$ of the left-hand side.
Additionally, the arguments of the functions are made implicit.

$$\langle\psi\rangle_v = \frac{1}{V}\int_{\mathcal V_f}\psi~\mathrm dV.$$

\subsection{The intrinsic phase average}

This average is defined as the mean value of some quantity $\psi$ over the fluid phase.

$$\langle\psi\rangle_f(\bm x, t) \doteq \frac{1}{V_f}\int_{\mathcal V} \psi(\bm x+\bm\xi, t) \delta(\bm x + \bm\xi, t) ~\left|\mathrm d\bm \xi\right|$$

Similarly to the fluid phase average, this formulation can be simplified.

$$\langle\psi\rangle_f = \frac{1}{V_f}\int_{\mathcal V_f}\psi~\mathrm dV$$

\subsection{Porosity}\label{porosity}

The porosity (volume ratio of fluid to solid phases) can be defined as the fluid phase average of the identity.

$$\epsilon \doteq \langle 1\rangle_v = \frac{V_f}{V}.$$

It is a link between both the averages used here.

\begin{equation}\label{eq:fluid-intrinsic}
  \langle\psi\rangle_v=\epsilon\langle\psi\rangle_f
\end{equation}

\section{Averaging rules}\label{averaging-rules}

To lighten up the derivation, we can define averaging rules.

\subsection{Linearity}\label{linearity}

The intrinsic phase average is a linear map. To prove this statement, let's average an affine function.

$$\langle a\psi + b\rangle_f = \frac{1}{V_f}\int_{\mathcal V_f}a\psi+b~\mathrm dV.$$

With constant coefficients $a$ and $b$, the quantity $\psi$ is the only remaining variable in the integral.

$$\langle a\psi + b\rangle_f = a \frac{1}{V_f}\int_{\mathcal V_f}\psi\mathrm dV + b.$$

Showing that the intrinsic phase average is a linear map.

$$\langle a\psi + b\rangle_f = a\langle\psi\rangle_f+b.$$

However, the fluid phase average is not.
It can be shown by developing the affine function.

$$\langle a\psi+b\rangle_v = \frac{1}{V}\int_{\mathcal V_f} a\psi+b~\mathrm dV.$$

Note that the normalizing volume $V$ is not equal to the integration domain's volume $\int_{\mathcal V_f}\mathrm dV=V_f$.

$$\langle a\psi+b\rangle_v = a\frac{1}{V}\int_{\mathcal V_f} \psi~\mathrm dV + b\frac{1}{V}\int_{\mathcal V_f}~\mathrm dV.$$

Thus, the fluid phase average of the constant coefficient $b$ is reduced by the porosity.

$$\langle a\psi+b\rangle_v = a\langle\psi\rangle_v+\epsilon b.$$

Showing that the fluid phase average is not a linear map.

\subsection{Integration of averaged variables}

Assuming that the control volume $\mathcal V$ is much larger than the porosity lenght-scale, the intrinsic average does not vary significantly within $\mathcal V$ \citep{Carbonell1984}.
Consequently, variables averaged over the intrinsic phase will be considered constant within an integrating (averaging) procedure over a subdomain $\mathcal D\subseteq\mathcal V$.

\begin{equation}\label{eq:integrate-average}
  \boxed{\int_{\mathcal D}\langle\psi\rangle_f~\mathrm dA\approx\langle\psi\rangle_f\int_{\mathcal D}~\mathrm dA.}
\end{equation}

$$\Rightarrow
\begin{cases}
    \langle\langle\psi\rangle_f\rangle_f=\langle\psi\rangle_f,\\
    \langle\langle\psi\rangle_f\rangle_v=\langle\psi\rangle_f\langle1\rangle_v=\epsilon\langle\psi\rangle_f.
\end{cases}$$

Note that it implies the same argument for the fluid phase average since the porosity $\epsilon=\langle1\rangle_v$ is an averaged quantity as well.

$$\epsilon\langle\psi\rangle_f\int_{\mathcal D}~\mathrm dA=\epsilon\int_{\mathcal D}\langle\psi\rangle_f~\mathrm dA,$$

$$\Leftrightarrow\epsilon\langle\psi\rangle_f\int_{\mathcal D}~\mathrm dA=\int_{\mathcal D}\epsilon\langle\psi\rangle_f~\mathrm dA,$$

$$\Leftrightarrow\langle\psi\rangle_v\int_{\mathcal D}~\mathrm dA=\int_{\mathcal D}\langle\psi\rangle_v~\mathrm dA,$$

\subsection{Nested averages}

Given that an averaged variable $\langle\psi\rangle_f$ should not vary significantly within the control volume $\mathcal V$, its own intrinsic average $\langle\langle\psi\rangle_f\rangle_f$ should be close to $\langle\psi\rangle_f$ itself.
This argument is consistent with the previous proposition (\autoref{eq:integrate-average}).

$$\langle\langle\psi\rangle_f\rangle_f=\langle\psi\rangle_f.$$

The averaged variable is thus not affected by a supplementary intrinsic phase average.
This rule is identical to the Reynolds averaging rule ${\bm{\overline{\bar u}}=\bm{\bar u}}$.
However, the fluid phase average will reduce the average value by the porosity (contrary to the previous assumption in literature ${\langle\langle\psi\rangle_v\rangle_v=\langle\psi\rangle_v}$ \citep{Gray1974}).

$$\langle\langle\psi\rangle_v\rangle_v = \epsilon\langle\langle\psi\rangle_v\rangle_f = \epsilon\langle\psi\rangle_v\left(=\epsilon^2\langle\psi\rangle_f\right),$$

$$\langle\langle\psi\rangle_v\rangle_f = \langle\epsilon\langle\psi\rangle_f\rangle_f = \epsilon\langle\langle\psi\rangle_f\rangle_f = \epsilon\langle\psi\rangle_f,$$

$$\langle\langle\psi\rangle_f\rangle_v = \epsilon\langle\langle\psi\rangle_f\rangle_f = \epsilon\langle\psi\rangle_f,$$

$$\langle\langle\psi\rangle_f\rangle_f = \langle\psi\rangle_f.$$

Note that the fluid and intrinsic phase average commute. This point is proof that the average of the spatial fluctuations is null.
Using the intrinsic phase average as a linear map and simplifying the nested average yields

$$\langle\psi\rangle_f = \langle\langle\psi\rangle_f + \tilde\psi\rangle_f,$$
$$\langle\psi\rangle_f = \langle\psi\rangle_f + \langle\tilde\psi\rangle_f.$$

And the average of the spatial fluctuations can only be null.

\begin{equation}\label{eq:average-fluctuation}
  \boxed{\langle\tilde\psi\rangle_f=0.}
\end{equation}

\subsection{Spatial averaging theorem}

Average and gradient do not commute as per the Leibniz formula.
By translating the control volume $\mathcal V$ through an arbitrary path, the directional derivative of an average can be related to the gradient of the averages.
An illustration of the procedure is given in \autoref{fig:trans}.

\begin{figure}
  \centering
  \includesvg{gfx/Whitaker_translation.svg}
  \caption{Differential displacement of the control volume \citep{Whitaker1999}.\label{fig:trans}}
\end{figure}

Starting with the spatial derivative in the form of a small translation, we come back to the definition of the derivative with a limit.

$$\frac{\mathrm d}{\mathrm d s}\int_{\mathcal V}\psi\delta~\mathrm dV = \lim_{\Delta s\to 0}\frac{1}{\Delta s}\left(\int_{\mathcal V(s+\Delta s)}\psi\delta~\mathrm dV - \int_{\mathcal V(s)}\psi\delta~\mathrm dV\right).$$

This expression is equivalent to taking the integral over the exclusive union of both control volumes ${\mathcal V(s+\Delta s)\oplus\mathcal V(s)}$.

$$\frac{\mathrm d}{\mathrm d s}\int_{\mathcal V}\psi\delta~\mathrm dV = \lim_{\Delta s\to 0}\frac{1}{\Delta s}\int_{\mathcal V(s+\Delta s)\oplus\mathcal V(s)}\psi\delta~\mathrm dV.$$

This integral can be simplified for an infinitesimal translation
$\Delta s$ \citep{SteinBarcellos1992}.

$$\frac{\mathrm d}{\mathrm d s}\int_{\mathcal V}\psi\delta~\mathrm dV = \lim_{\Delta s\to 0}\frac{1}{\Delta s}\int_{\mathcal S}\psi\delta\bm\lambda\cdot\bm n~\mathrm dA\cdot \Delta s.$$

The limit on $\Delta s$ vanishes and the constant vector $\bm\lambda$ can be taken out of the integral.

$$\frac{\mathrm d}{\mathrm d s}\int_{\mathcal V}\psi\delta~\mathrm dV = \bm\lambda\cdot\int_{\mathcal S}\psi\delta\bm n~\mathrm dA.$$

The left hand side can be rewritten in terms of $\bm\lambda$ by taking the directional derivative.

$$\bm\lambda\cdot\nabla\left(\int_{\mathcal V}\psi\delta~\mathrm dV\right) = \bm\lambda\cdot\int_{\mathcal S}\psi\bm n~\mathrm dA.$$

Since the above equation holds for any translation $\bm\lambda\in\mathbb R^3$, we have that

$$\nabla\left(\int_{\mathcal V}\psi\delta~\mathrm dV\right) = \int_{\mathcal S}\psi\delta\bm n~\mathrm dA.$$

By subtracting the fluid-solid interface $\mathcal S_s$ from the control volume's border $\partial\mathcal V$, the indicating function can be disposed of in the previous expression.
This can be seen in \autoref{fig:trans} by keeping in mind that the indicating function is null outside of the fluid phase's domain $\mathcal V_f$.

$$\nabla\left(\int_{\mathcal V}\psi\delta~\mathrm dV\right) = \int_{\partial\mathcal V_f}\psi\bm n~\mathrm dA - \int_{\mathcal S_s}\psi\bm n~\mathrm dA.$$

This form allows the application of the divergence theorem.

$$\nabla\left(\int_{\mathcal V}\psi\delta~\mathrm dV\right) = \int_{\mathcal V_f}\nabla\psi~\mathrm dV - \int_{\mathcal S_s}\psi\bm n~\mathrm dA.$$

And because the control volume $V$ is fixed, the fluid phase average smoothly appears.

$$\nabla\left(\frac{1}{V}\int_{\mathcal V}\psi\delta~\mathrm dV\right) = \frac{1}{V}\int_{\mathcal V_f}\nabla\psi~\mathrm dV - \frac{1}{V}\int_{\mathcal S_s}\psi\bm n~\mathrm dA.$$

The spatial averaging theorem finally takes the following form :

\begin{equation}\label{eq:averaging-theorem}
  \boxed{\langle\nabla\psi\rangle_v = \nabla\langle\psi\rangle_v + \frac{1}{V}\int_{\mathcal S_s}\psi\bm n~\mathrm dA.}
\end{equation}

Note that this theorem is valid not only for gradients of scalars ${\nabla\psi}$ but for divergences of vectors ${\bm\nabla\cdot\bm\psi}$ and jacobians ${\nabla\bm\psi}$ as well.

$$\langle\nabla\cdot\bm\psi\rangle_v = \nabla\cdot\langle\bm\psi\rangle_v + \frac{1}{V}\int_{\mathcal S_s}\bm\psi\cdot\bm n~\mathrm dA,$$

$$\langle\nabla\bm\psi\rangle_v = \nabla\langle\bm\psi\rangle_v + \frac{1}{V}\int_{\mathcal S_s}\bm\psi\otimes\bm n~\mathrm dA.$$

\subsection{Spatially averaging a constant variable}

The spatial averaging theorem applied to a constant variable expresses the changes in porosity.

$$\langle\nabla 1\rangle_v = \nabla\langle1\rangle_v + \frac{1}{V}\int_{\mathcal S_s}\bm n~\mathrm dA,$$
$$\langle 0\rangle_v = \nabla\epsilon + \frac{1}{V}\int_{\mathcal S_s}\bm n~\mathrm dA,$$

\begin{equation}\label{eq:average-empty}
  \boxed{\frac{1}{V}\int_{\mathcal S_s}\bm n~\mathrm dA = -\nabla\epsilon.}
\end{equation}

\subsection{Spatially averaging a time derivative}

It can be demonstrated that for fixed control volumes (with a null velocity at the boundary, ${\bm w=0}\ \forall\bm x\in\mathcal S_s$), the average of the accumulation does not include any additional terms.

$$\frac{\partial}{\partial t}\int_{\mathcal V_f}\psi~\mathrm dV = \int_{\mathcal V_f}\frac{\partial\psi}{\partial t}~\mathrm dV + \int_{\partial\mathcal V_f}\psi \cdot(\bm w\cdot\bm n)~\mathrm dA,$$

\begin{equation}\label{eq:average-accumulation}
  \boxed{\frac{\partial}{\partial t}\int_{\mathcal V_f}\psi~\mathrm dV = \int_{\mathcal V_f}\frac{\partial\psi}{\partial t}~\mathrm dV.}
\end{equation}

\subsection{No-slip condition}

At the interface between the fluid and the solid phases $\mathcal S_s$, the velocity of the fluid phase $\bm u$ and that of the solid phase $\bm w$ are one and the same. Because the scope of this derivation is limited to fixed beds, the solid phase's velocity $\bm w$ is null.

$$\bm u = \bm w = \bm 0,\ \forall\bm x\in\mathcal S_s.$$

Consequently, the temporal average of the velocity at that point is null as well.

$$\bm u = \bm{\bar u} = \bm 0,\ \forall\bm x\in\mathcal S_s.$$

It also implies null temporal fluctuations $\bm u'$ are null since $\bm u=\bm{\bar u}+\bm u'$.

\begin{equation}\label{eq:no-slip}
  \bm u = \bm{\bar u} = \bm u' = \bm 0,\ \forall\bm x\in\mathcal S_s.
\end{equation}

This relation will greatly simplify the expressions to come.

\section{Derivation of the DANS equations}

We apply the fluid phase average to the Reynolds-(or time-)averaged Navier-Stokes equation, aiming to develop a new set of equations with the double-averaged variables as the unknowns.

\subsection{Continuity equation}

The continuity equation is fundamentally unchanged. The original continuity equation is

$$\langle\bm\nabla\cdot\bm{\bar u}\rangle_v = 0.$$

By applying the spatial averaging theorem (\autoref{eq:averaging-theorem}), we have that

$$\bm\nabla\cdot\langle\bm{\bar u}\rangle_v + \frac{1}{V}\int_{\mathcal S_s}{\bm{\bar u}\cdot \bm n}~\mathrm dA = 0.$$

Because of the no-slip condition (\autoref{eq:no-slip}) at the fluid-solid interface, both the time-averaged speed ${\bm{\bar u}}$ and the temporal fluctuations ${\bm u'}$ are null as well.

$$\bm\nabla\cdot\langle\bm{\bar u}\rangle_v = 0.$$

By expressing the intrinsic phase average (\autoref{eq:fluid-intrinsic}), the double-averaged continuity equation is

$$\boxed{\bm\nabla\cdot\left(\epsilon\langle\bm{\bar u}\rangle_f\right) = 0.}$$

\subsection{Momentum equation}

\subsubsection{Advective term}

The advective term is the source of turbulence. When its velocity is decomposed (\autoref{eq:decomposition}), the dispersive stress appears.

$$\langle\bm\nabla\cdot \bm{\bar u}\otimes\bm{\bar u}\rangle_v = \nabla\cdot\langle\bm{\bar u}\otimes\bm{\bar u}\rangle_v + \int_{\mathcal S_s}\bm{\bar u}\otimes\bm{\bar u}\cdot\bm n~\mathrm dA.$$

But the integral term is null because of the no-slip condition
(\autoref{eq:no-slip}).

$$\langle\bm\nabla\cdot \bm{\bar u}\otimes\bm{\bar u}\rangle_v = \nabla\cdot\langle\bm{\bar u}\otimes\bm{\bar u}\rangle_v.$$

We can take the intrinsic phase average (\autoref{eq:fluid-intrinsic}).

$$\langle\bm\nabla\cdot \bm{\bar u}\otimes\bm{\bar u}\rangle_v = \nabla\cdot\left(\epsilon\langle\bm{\bar u}\otimes\bm{\bar u}\rangle_f\right).$$

And decompose the speed (\autoref{eq:decomposition}) to get the dispersive stress tensor.
Remember that the average of the spatial fluctuations is null (\autoref{eq:average-fluctuation}) and that averaged quantities can be considered constant within an averaging procedure (\autoref{eq:integrate-average}).

$$\langle\bm\nabla\cdot \bm{\bar u}\otimes\bm{\bar u}\rangle_v = \nabla\cdot\left(\epsilon\langle\bm{\bar u}\rangle_f\otimes\langle\bm{\bar u}\rangle_f + \epsilon\langle \tilde{\bm u}\otimes\tilde{\bm u}\rangle_f\right).$$

Distributing the divergence yields the final expression.

$$\boxed{\langle\bm\nabla\cdot \bm{\bar u}\otimes\bm{\bar u}\rangle_v = \nabla\cdot\left(\epsilon\langle\bm{\bar u}\rangle_f\otimes\langle\bm{\bar u}\rangle_f\right) + \nabla\cdot\left(\epsilon\langle \tilde{\bm u}\otimes\tilde{\bm u}\rangle_f\right).}$$

The same procedure (until the decomposition $\bm{\bar u}=\bm u' + \tilde{\bm u}$) can be followed for the Reynolds stress tensor.

$$\boxed{\langle\bm\nabla\cdot \overline{\bm u'\otimes\bm u'}\rangle_v = \nabla\cdot\left(\epsilon\langle\overline{\bm u'\otimes\bm u'}\rangle_f\right).}$$

And the total advection is as follows.

$$\boxed{\varrho\langle\overline{\bm\nabla\cdot \bm u\otimes\bm u}\rangle_v =
\underbrace{\varrho\nabla\cdot\left(\epsilon\langle\bm{\bar u}\rangle_f\otimes\langle\bm{\bar u}\rangle_f\right)}_{\text{Double-averaged advection}}
+ \underbrace{\varrho\nabla\cdot\left(\epsilon\langle \tilde{\bm u}\otimes\tilde{\bm u}\rangle_f\right)}_{\text{Dispersive stress}}
+ \underbrace{\varrho\nabla\cdot\left(\epsilon\langle\overline{\bm u'\otimes\bm u'}\rangle_f\right)}_{\text{Reynolds stress}}.}$$

\subsubsection{Accumulative term}

From the average of a time derivative (\autoref{eq:average-accumulation}), we know that the average and the derivative commute for a fixed control volume ($\bm w = \bm 0$).

$$\left\langle\frac{\partial \bm{\bar u}}{\partial t}\right\rangle_v = \epsilon\left\langle\frac{\partial \bm{\bar u}}{\partial t}\right\rangle_f = \epsilon\frac{\partial\langle\bm{\bar u}\rangle_f}{\partial t}.$$

\subsubsection{Pressure gradient}

The pressure is subject to sudden spatial fluctuations close to rigid walls.
Thus the decomposition shows a integral term along the fluid-solid interface expressing part of the drag force density.
We can start by applying the spatial averaging theorem (\autoref{eq:averaging-theorem}).

$$\langle\nabla\bar p\rangle_v = \nabla\langle\bar p\rangle_v + \frac{1}{V}\int_{\mathcal S_s}\bar p\bm n~\mathrm dA.$$

The intrinsic phase pressure can directly be expressed (\autoref{eq:decomposition}),

$$\langle\nabla\bar p\rangle_v = \nabla \left(\epsilon\langle\bar p\rangle_f\right) + \frac{1}{V}\int_{\mathcal S_s}\left(\langle\bar p\rangle_f+\tilde p\right)\bm n~\mathrm dA,$$

the averaged pressure can be taken out of the integral (\autoref{eq:integrate-average}), and the gradient distributed.

$$\langle\nabla\bar p\rangle_v = \langle\bar p\rangle_f \nabla \epsilon + \epsilon\nabla\langle\bar p\rangle_f + \langle\bar p\rangle_f\frac{1}{V}\int_{\mathcal S_s} \bm n~\mathrm dA + \frac{1}{V}\int_{\mathcal S_s}\tilde p\bm n~\mathrm dA.$$

By means of the spatial averaging theorem applied to a constant variable (\autoref{eq:average-empty}), the empty integral is simplified.

$$\langle\nabla\bar p\rangle_v = \langle\bar p\rangle_f \nabla \epsilon + \epsilon\nabla\langle\bar p\rangle_f - \langle\bar p\rangle_f\nabla\epsilon + \frac{1}{V}\int_{\mathcal S_s}\tilde p\bm n~\mathrm dA.$$

And the final expression is obtained.

$$\boxed{\langle\nabla\bar p\rangle_v = \epsilon \nabla\langle\bar p\rangle_f + \underbrace{\frac{1}{V}\int_{\mathcal S_s}\tilde p\bm n~\mathrm dA}_{\text{Contribution to the drag force density }\bm{\bar f}}.}$$

\subsubsection{Diffusion}

The diffusion term receives a correction due to porosity changes (${\nabla\epsilon}$ and ${\nabla^2\epsilon}$).
Additionally, the contribution of spatial fluctuations of velocity to the drag force density across the fluid-solid interface appears.
First of all, the spatial averaging theorem can be applied twice (\autoref{eq:averaging-theorem}).

$$\langle\bm\nabla^2\bm{\bar u}\rangle_v = \langle\bm\nabla\cdot\nabla\bm{\bar u}\rangle_v,$$
$$\langle\bm\nabla^2\bm{\bar u}\rangle_v = \bm\nabla\cdot\langle\nabla\bm{\bar u}\rangle_v + \frac{1}{V}\int_{\mathcal S_s}\nabla\bm{\bar u}\cdot \bm n~\mathrm dA,$$
$$\langle\bm\nabla^2\bm{\bar u}\rangle_v = \bm\nabla\cdot\left(\nabla\langle\bm{\bar u}\rangle_v+\frac{1}{V}\int_{\mathcal S_s}\bm{\bar u}\otimes\bm n~\mathrm dA\right) + \frac{1}{V}\int_{\mathcal S_s}\nabla\bm{\bar u}\cdot \bm n~\mathrm dA.$$

The no-slip condition is applied (\autoref{eq:no-slip}).

$$\langle\bm\nabla^2\bm{\bar u}\rangle_v = \bm\nabla\cdot\nabla\langle\bm{\bar u}\rangle_v + \frac{1}{V}\int_{\mathcal S_s}\nabla\bm{\bar u}\cdot \bm n~\mathrm dA.$$

This last term is an expression of the drag forces applied by the solid boundaries.
While closures exist for the drag force density for homogeneous media, we wish to extract the effects of heterogeneity $\left(\nabla\epsilon\right)$ from the integral term and then use an existing closure for the homogeneous component.

\bigskip

The intrinsic phase average (\autoref{eq:fluid-intrinsic}) can be expressed to decompose the velocity vector (\autoref{eq:decomposition}).

$$\langle\bm\nabla^2\bm{\bar u}\rangle_v = \bm\nabla\cdot\nabla\left(\epsilon\langle\bm{\bar u}\rangle_f\right) + \frac{1}{V}\int_{\mathcal S_s}\left(\nabla\langle \bm{\bar u}\rangle_f\cdot \bm n + \nabla\tilde{\bm u}\cdot \bm n\right)~\mathrm dA,$$

$$\langle\bm\nabla^2\bm{\bar u}\rangle_v = \bm\nabla\cdot\nabla\left(\epsilon\langle\bm{\bar u}\rangle_f\right) + \frac{1}{V}\int_{\mathcal S_s}\nabla\langle \bm{\bar u}\rangle_f\cdot \bm n~\mathrm dA + \frac{1}{V}\int_{\mathcal S_s}\nabla\tilde{\bm u}\cdot \bm n~\mathrm dA,$$

$$\langle\bm\nabla^2\bm{\bar u}\rangle_v = \epsilon\nabla^2\langle\bm{\bar u}\rangle_f + 2\nabla\langle\bm{\bar u}\rangle_f\cdot\nabla\epsilon + \langle\bm{\bar u}\rangle_f\nabla^2\epsilon + \frac{1}{V}\int_{\mathcal S_s}\nabla\langle \bm{\bar u}\rangle_f\cdot \bm n~\mathrm dA + \frac{1}{V}\int_{\mathcal S_s}\nabla\tilde{\bm u}\cdot \bm n~\mathrm dA.$$

While \citet{Rousseau2022} directly wrote that ${\int_\mathcal{S_s}\nabla\langle\bm{\bar u}\rangle_f\cdot\bm n~\mathrm dA=\nabla\langle\bm{\bar u}\rangle_f\cdot\int_\mathcal{S_s}\bm n~\mathrm dA}$, it can be quickly justified by switching between fluid and intrinsic phase averages (\autoref{eq:fluid-intrinsic}), exposing the average of the gradient (\autoref{eq:averaging-theorem}) to simplify the integral (\autoref{eq:integrate-average}), and using the no-slip condition (\autoref{eq:no-slip}).

\begin{align*}
\int_\mathcal{S_s}\nabla\langle\bm{\bar u}\rangle_f\cdot\bm n~\mathrm dA &= \int_{\mathcal S_s}\nabla\left(\epsilon^{-1}\langle\bm{\bar u}\rangle_v\right)\cdot\bm n~\mathrm dA,\\
    &= \int_{\mathcal S_s}\nabla\langle\epsilon^{-1}\bm{\bar u}\rangle_v\cdot\bm n~\mathrm dA,\\
    &= \int_{\mathcal S_s}\left[\left\langle\nabla\left(\epsilon^{-1}\bm{\bar u}\right)\right\rangle_v - \int_{\mathcal S_s}\epsilon^{-1}\bm{\bar u}\otimes\bm n~\mathrm dA\right]\cdot\bm n~\mathrm dA,\\
    &= \int_{\mathcal S_s}\left\langle\nabla\left(\epsilon^{-1}\bm{\bar u}\right)\right\rangle_v\cdot\bm n~\mathrm dA,\\
    &= \left\langle\nabla\left(\epsilon^{-1}\bm{\bar u}\right)\right\rangle_v\cdot\int_{\mathcal S_s}\bm n~\mathrm dA,\\
    &= -\left\langle\nabla\left(\epsilon^{-1}\bm{\bar u}\right)\right\rangle_v\cdot\nabla\epsilon,\\
    &= -\left(\nabla\left\langle\epsilon^{-1}\bm{\bar u}\right\rangle_v + \int_{\mathcal S_s}\epsilon^{-1}\bm{\bar u}\otimes\bm n~\mathrm dA\right)\cdot\nabla\epsilon,\\
    &= -\nabla\left\langle\epsilon^{-1}\bm{\bar u}\right\rangle_v\cdot\nabla\epsilon,\\
    &= -\nabla\left(\epsilon^{-1}\left\langle\bm{\bar u}\right\rangle_v\right)\cdot\nabla\epsilon,\\
    &= -\nabla\left\langle\bm{\bar u}\right\rangle_f\cdot\nabla\epsilon.\\
\end{align*}

Inserting this result back into the laplacian of the velocity vector ${\langle\nabla^2\bm{\bar u}\rangle_v}$ yields

$$\langle\bm\nabla^2\bm{\bar u}\rangle_v = \bm\nabla\cdot\nabla\left(\epsilon\langle\bm{\bar u}\rangle_f\right) -\nabla\langle \bm{\bar u}\rangle_f\nabla\epsilon + \frac{1}{V}\int_{\mathcal S_s}\nabla\tilde{\bm u}\cdot \bm n~\mathrm dA.$$

The divergence can be developed to read the total effect of a heterogeneous porosity. When terms involving the porosity gradient ${\nabla\epsilon}$ and laplacian ${\nabla^2\epsilon}$ are omitted (homogeneous medium), the remaining laplacian of the average velocity is termed the \emph{Brinkman correction}.

$$\langle\bm\nabla^2\bm{\bar u}\rangle_v = \bm\nabla\cdot\left[\epsilon\nabla\langle\bm{\bar u}\rangle_f + \langle\bm{\bar u}\rangle_f\otimes\nabla\epsilon\right] - \nabla\langle \bm{\bar u}\rangle_f\nabla\epsilon + \frac{1}{V}\int_{\mathcal S_s}\nabla\tilde{\bm u}\cdot \bm n~\mathrm dA.$$

$$\langle\bm\nabla^2\bm{\bar u}\rangle_v = \underbrace{\epsilon\nabla^2\langle\bm{\bar u}\rangle_f}_\text{Brinkman correction} + 2\nabla\langle\bm{\bar u}\rangle_f\cdot\nabla\epsilon + \langle\bm{\bar u}\rangle_f\nabla^2\epsilon - \nabla\langle \bm{\bar u}\rangle_f\nabla\epsilon + \frac{1}{V}\int_{\mathcal S_s}\nabla\tilde{\bm u}\cdot \bm n~\mathrm dA.$$

Which leads to the final expression.

$$\boxed{\langle\bm\nabla^2\bm{\bar u}\rangle_v = \underbrace{\epsilon\nabla^2\langle\bm{\bar u}\rangle_f}_\text{Brinkman correction} + \underbrace{\nabla\langle\bm{\bar u}\rangle_f\cdot\nabla\epsilon + \langle\bm{\bar u}\rangle_f\nabla^2\epsilon}_{\text{Porosity heterogeneity}} + \underbrace{\frac{1}{V}\int_{\mathcal S_s}\nabla\tilde{\bm u}\cdot \bm n~\mathrm dA}_{\text{Contribution to the drag force density }\bm{\bar f}}.}$$

\subsubsection{Drag force density}

The contribution of the fluctuations of the pressure and of the laplacian can be combined into a single term.

$$\bm{\bar f} \doteq \frac{1}{V_f}\int_{\mathcal S_s}\left(\mu\nabla\tilde{\bm u} - \tilde p\bm I\right)\bm n~\mathrm dA.$$

By taking a step back, it can be noted that these are akin to the fluctuations of the applied traction ${\bm t=\bm\sigma\bm n}$ (where ${\bm\sigma}$ is the stress tensor) on the solid interface.

$$\bm{\bar f} = \frac{1}{V_f}\int_{\mathcal S_s}\underbrace{\tilde{\bm\sigma}\bm n}_{\tilde{\bm t}}~\mathrm dA,$$

$$\bm\sigma = \mu\nabla\bm u - p.$$

\paragraph{Closure of the drag force density}

To have a second-order closure, we can use Ergun's equation and take the Forchheimer
parametrisation \citep{W1996}:

$$\bm{ f} = \overbrace{-\underbrace{\frac{\epsilon\mu}{K}\langle\bm u\rangle_f}_{\text{Darcy}} - \underbrace{\frac{\varrho\epsilon^2c}{\sqrt{K}}\langle\bm u\rangle_f\left|\langle\bm u\rangle_f\right|}_{\text{Forchheimer correction}}}^{\text{Ergun's drag force density}}.$$

which is already a space-averaged form. Time-averaging this form and decomposing the velocity ${\bm u = \bm{\bar u} + \bm u'}$ yields \citep{G2000}:

$$\bm{\bar f} = -\frac{\epsilon\mu}{K}\langle\bm{\bar u}\rangle_f - \frac{\varrho\epsilon^2c}{\sqrt{K}}\overline{\langle\bm u\rangle_f\left|\langle\bm u\rangle_f\right|},$$

$$\bm{\bar f} = -\frac{\epsilon\mu}{K}\langle\bm{\bar u}\rangle_f - \frac{\varrho\epsilon^2c}{\sqrt{K}}\overline{\left(\langle\bm{\bar u}\rangle_f + \langle\bm u'\rangle_f\right)\left|\langle\bm{\bar u}\rangle_f + \langle\bm u'\rangle_f\right|}.$$

Assuming the temporal fluctuations $\bm u'$ are much smaller than the temporally averaged values $\bm{\bar u}$ in the calculation of the norm, it can be expanded to the first-order as follows.
Note that the assumption $\bm u'\ll\bm{\bar u}$ can be made because the drag force density is predominant (in the momentum balance equation) when the flow is laminar (weakly turbulent or not at all).

\begin{align*}
\left|\langle\bm{\bar u}\rangle_f+\langle \bm u'\rangle_f\right| & = \sqrt{\left|\langle\bm{\bar u}\rangle_f\right|^2+2\langle\bm{\bar u}\rangle_f\cdot\langle\bm u'\rangle_f+\left|\langle \bm u'\rangle_f\right|^2},\\
    & = \sqrt{\left|\langle\bm{\bar u}\rangle_f\right|^2+2\langle\bm{\bar u}\rangle_f\cdot\langle\bm u'\rangle_f},\\
    &= \left|\langle\bm{\bar u}\rangle_f\right|+\frac{\langle\bm{\bar u}\rangle_f\cdot\langle\bm u'\rangle_f}{\left|\langle\bm{\bar u}\rangle_f\right|} + \mathcal \varepsilon\left(\left|\langle\bm u'\rangle_f\right|\right).
\end{align*}

Reinjecting this expansion in the time-averaged drag force density, we get that

\begin{align*}
\bm{\bar f} &= -\frac{\epsilon\mu}{K}\langle\bm{\bar u}\rangle_f - \frac{\varrho\epsilon^2c}{\sqrt{K}}\overline{\left(\langle\bm{\bar u}\rangle_f + \langle\bm u'\rangle_f\right)\left(\left|\langle\bm{\bar u}\rangle_f\right|+\frac{\langle\bm{\bar u}\rangle_f\cdot\langle\bm u'\rangle_f}{\left|\langle\bm{\bar u}\rangle_f\right|}\right)},\\
    &= -\frac{\epsilon\mu}{K}\langle\bm{\bar u}\rangle_f - \frac{\varrho\epsilon^2c}{\sqrt{K}}\left(\overline{\langle\bm{\bar u}\rangle_f\left|\langle\bm{\bar u}\rangle_f\right| + \langle\bm u'\rangle_f\left|\langle\bm{\bar u}\rangle_f\right| + \langle\bm{\bar u}\rangle_f\frac{\langle\bm{\bar u}\rangle_f\cdot\langle\bm u'\rangle_f}{\left|\langle\bm{\bar u}\rangle_f\right|} + \langle\bm u'\rangle_f\frac{\langle\bm{\bar u}\rangle_f\cdot\langle\bm u'\rangle_f}{\left|\langle\bm{\bar u}\rangle_f\right|}}\right),\\
    &= -\frac{\epsilon\mu}{K}\langle\bm{\bar u}\rangle_f - \frac{\varrho\epsilon^2c}{\sqrt{K}}\left(\langle\bm{\bar u}\rangle_f\left|\langle\bm{\bar u}\rangle_f\right| + \overline{\langle\bm u'\rangle_f\frac{\langle\bm{\bar u}\rangle_f\cdot\langle\bm u'\rangle_f}{\left|\langle\bm{\bar u}\rangle_f\right|}}\right),\\
    &= -\frac{\epsilon\mu}{K}\langle\bm{\bar u}\rangle_f - \frac{\varrho\epsilon^2c}{\sqrt{K}}\left(\langle\bm{\bar u}\rangle_f\left|\langle\bm{\bar u}\rangle_f\right| + \overline{\langle\bm u'\rangle_f\otimes\langle\bm u'\rangle_f}\frac{\langle\bm{\bar u}\rangle_f}{\left|\langle\bm{\bar u}\rangle_f\right|}\right).
\end{align*}

Which concludes the derivation led by \citet{G2000} and \citet{Kuwata2013} to time-average Ergun's drag force density.

$$\boxed{\bm{\bar f} = -\frac{\epsilon\mu}{K}\langle\bm{\bar u}\rangle_f - \frac{\varrho\epsilon^2c}{\sqrt{K}}\left(\langle\bm{\bar u}\rangle_f\left|\langle\bm{\bar u}\rangle_f\right| + \overline{\langle\bm u'\rangle_f\otimes\langle\bm u'\rangle_f}\frac{\langle\bm{\bar u}\rangle_f}{\left|\langle\bm{\bar u}\rangle_f\right|}\right).}$$

\subsection{Final system of equations}

The new set of equations, the double-averaged Navier-Stokes equations, is as follows.

\begin{enumerate}
\def\labelenumi{\arabic{enumi}.}
\item
  Continuity equation
\end{enumerate}

$$\bm\nabla\cdot\left(\epsilon\langle\bm{\bar u}\rangle_f\right) = 0.$$

\begin{enumerate}
\def\labelenumi{\arabic{enumi}.}
\setcounter{enumi}{1}
\item
  Momentum equation
\end{enumerate}

\begin{align*}
    \epsilon\frac{\partial\langle\bm{\bar u}\rangle_f}{\partial t} + \bm\nabla\cdot\left(\epsilon\langle\bm{\bar u}\rangle_f\otimes\langle\bm{\bar u}\rangle_f\right) =&+ \epsilon\bm g\\
        &- \frac{\epsilon}{\varrho}\nabla\langle\overline p\rangle_f\\
        &+ \bm\nabla\cdot\left(\nu\nabla\cdot\left(\epsilon\langle\bm{\bar u}\rangle_f\right)\right) -\nu\nabla\langle\bm{\bar u}\rangle_f\nabla\epsilon\\
        &- \bm\nabla\cdot\left(\epsilon\langle \overline{\bm u'\otimes\bm u'}\rangle_f\right)\\
        &- \bm\nabla\cdot\left(\epsilon\langle \tilde{\bm u}\otimes\tilde{\bm u}\rangle_f\right)\\
        &+ \frac{\epsilon}{\varrho}\bm{\bar f},\\
\end{align*}

Where the last three terms need closure.

\clearpage

\bibliography{references}
\label{sec:bibliographie}
\addcontentsline{toc}{section}{\bibname}
\thispagestyle{empty}

\end{document}